\documentclass{article}
\usepackage{preamble}

\definecolor{cluster1}{RGB}{141,211,199}
\definecolor{cluster2}{RGB}{255,255,179}
\definecolor{cluster3}{RGB}{190,186,218}
\definecolor{cluster4}{RGB}{251,128,114}

\begin{document}
\header{Filemon Mateus}{CS 6140}{Clustering}{Homework \# 4}{\today}
\begin{enumerate}[leftmargin={*}, font={\bf}, label={\arabic*.}, ref={\arabic*}]
  \item \label{qst:1}
    \begin{enumerate}[label={(\Alph*)}, ref={\theenumi(\Alph*)}]
      \item \label{qst:1A}
        Below is a tabulated list of sets corresponding to the resulting clusters of
        running a fully-fledge vanilla implementation of {\sf Agglomerative Clustering}
        on {\sf C1.txt} dataset.
        \begin{table}[H]
          \centering
          \begin{tabular}{|l|c|c|c|}
            \hline
            & {\bf Cluster 1} & {\bf Cluster 2} & {\bf Cluster 3} \\ \hline
            {\bf Single-Link} & $\{1,2\}$ & $\{3,4\}$ & $\{5 \cdots 21\}$ \\
            {\bf Complete-Link} & $\{1,2\}$ & $\{3,4,5,6,7,15,16,17,18,19,20,21\}$ & $\{8 \cdots 14\}$ \\ \hline
          \end{tabular}
          \caption{
            Hierarchical Agglomerative Clustering. See \autoref{lst:hierarchical} \\
            for a fuller explanation on how these clusters have been obtained.
          }
          \label{tbl:hierarchical}
        \end{table}

        Pictorially, \autoref{tbl:hierarchical} translates to the following plots:
        \begin{figure}[H]
          \centering
          \begin{subfigure}[b]{0.43\linewidth}
            \begin{tikzpicture}
              \begin{axis}[
                width={\linewidth},
                legend pos={north east},
                only marks,
                mark size=3pt,
                cycle list={
                  {draw=red!90!black, mark=+},
                  {draw=blue!90!black, mark=-},
                  {draw=green!60!black, mark=x}
                }
              ]
                \pgfplotsinvokeforeach{1,2,3}{
                  \addplot+ table[col sep=comma] {dev-dir/week4/data/hierarchical/s-link-cluster-#1.csv};
                  \addlegendentry{Cluster #1};
                }
              \end{axis}
            \end{tikzpicture}
            \caption{Single Linkage Hierarchical Clustering.}
          \end{subfigure}
          \begin{subfigure}[b]{0.43\linewidth}
            \begin{tikzpicture}
              \begin{axis}[
                width={\linewidth},
                legend pos={north east},
                only marks,
                mark size=3pt,
                cycle list={
                  {draw=red!90!black, mark=+},
                  {draw=blue!90!black, mark=-},
                  {draw=green!60!black, mark=x}
                }
              ]
                \pgfplotsinvokeforeach{1,2,3}{
                  \addplot+ table[col sep=comma] {dev-dir/week4/data/hierarchical/c-link-cluster-#1.csv};
                  \addlegendentry{Cluster #1};
                }
              \end{axis}
            \end{tikzpicture}
            \caption{Complete Linkage Hierarchical Clustering.}
          \end{subfigure}
        \end{figure}

      \item \label{qst:1B}
        Below is an implementation of {\sf AgglomerativeClustering} in {\sf scikit-learn}:
        \lstinputlisting
          [caption={\sf hierarchical-scikit.py}, label={lst:hierarchical-scikit}]
          {dev-dir/week4/hierarchical-scikit.py}

      \item \label{qst:1C}
        The corresponding dendrograms are plotted below.
        \begin{figure}[H]
          \centering
          \begin{subfigure}[b]{0.43\linewidth}
            \includegraphics[scale=0.4]{dev-dir/week4/figs/dendrogram-single.pdf}
            \caption{Single Linkage Hierarchical Clustering.}
          \end{subfigure}
          \begin{subfigure}[b]{0.43\linewidth}
            \includegraphics[scale=0.4]{dev-dir/week4/figs/dendrogram-complete.pdf}
            \caption{Complete Linkage Hierarchical Clustering.}
          \end{subfigure}
          \caption{
            Hierarchical Clustering Dendrograms. See \autoref{lst:hierarchical-dendrogram} \\
            for a fuller explanation on how these dendrograms have been generated.
          }
        \end{figure}
    \end{enumerate}

  \item \label{qst:2}
    \begin{enumerate}[label={(\Alph*)}, ref={\theenumi(\Alph*)}]
      \item \label{qst:2A}
        \begin{enumerate}[label={(\roman*)}, ref={\theenumii(\roman*)}]
          \item
            Below is a plot corresponding to the resulting output of a fully-fledge vanilla
            implementation of {\sf Gonzalez} and {\sf Kmeans{\tt++}} clustering algorithms
            on {\sf C2.txt} dataset.
            \begin{figure}[H]
              \centering
              \begin{subfigure}[b]{0.43\linewidth}
                \begin{tikzpicture}
                  \begin{axis}[
                    width={\linewidth},
                    legend pos={north east},
                    legend entries={Cluster 1, Cluster 2, Cluster 3, Cluster 4},
                    point meta={explicit},
                    mark size={1.5pt},
                    only marks,
                  ]
                    \addplot [
                      scatter,
                      scatter/classes={
                        0={fill={cluster1}, draw={cluster1!80!black}},
                        1={fill={cluster2}, draw={cluster2!80!black}},
                        2={fill={cluster3}, draw={cluster3!80!black}},
                        3={fill={cluster4}, draw={cluster4!80!black}}
                      },
                    ] table[col sep=comma, x index=0, y index=1, meta index=2] {dev-dir/week4/data/gonzalez/points.csv};
                    \addplot [
                      mark=x,
                      mark size=5pt,
                      color=red
                    ] table[col sep=comma] {dev-dir/week4/data/gonzalez/centers.csv};
                  \end{axis}
                \end{tikzpicture}
                \caption{{\sf Gonzalez} Clustering Algorithm.}
              \end{subfigure}
              \begin{subfigure}[b]{0.43\linewidth}
                \begin{tikzpicture}
                  \begin{axis}[
                    width={\linewidth},
                    legend pos={north east},
                    legend entries={Cluster 1, Cluster 2, Cluster 3, Cluster 4},
                    point meta={explicit},
                    mark size={1.5pt},
                    only marks,
                  ]
                    \addplot [
                      scatter,
                      scatter/classes={
                        0={fill={cluster1}, draw={cluster1!80!black}},
                        1={fill={cluster2}, draw={cluster2!80!black}},
                        2={fill={cluster3}, draw={cluster3!80!black}},
                        3={fill={cluster4}, draw={cluster4!80!black}}
                      },
                    ] table[col sep=comma, x index=0, y index=1, meta index=2] {dev-dir/week4/data/kplusplus/points.csv};
                    \addplot [
                      mark=x,
                      mark size=5pt,
                      color=red
                    ] table[col sep=comma] {dev-dir/week4/data/kplusplus/centers.csv};
                  \end{axis}
                \end{tikzpicture}
                \caption{{\sf Kmeans{\tt++}} Clustering Algorithm.}
              \end{subfigure}
              \caption{
                {\sf Gonzalez} and {\sf Kmeans{\tt++}} clustering algorithms on {\sf C2.txt}
                dataset. Please refer to \\ \autoref{lst:gonzalez} and \ref{lst:kplusplus} for
                a fuller explanation on how these figures have been generated.
              }
            \end{figure}

          \item
            The following $4$-center costs were obtained during the run:
            \begin{table}[H]
              \centering
              \begin{tabular}{|c|c|c|}
                \hline
                & {\sf Gonzalez} & {\sf Kmeans{\tt++}} \\ \hline
                {\sf 4-Center Cost} & {\bf 3.481123664686921} & {\bf 11.348260306157117} \\ \hline
              \end{tabular}
              \caption{4-center cost.}
            \end{table}

          \item
            The following $4$-means costs were obtained during the run:
            \begin{table}[H]
              \centering
              \begin{tabular}{|c|c|c|}
                \hline
                & {\sf Gonzalez} & {\sf Kmeans{\tt++}} \\ \hline
                {\sf 4-Means Cost} & {\bf 1.6005365908052087} & {\bf 1.0416802543222048} \\ \hline
              \end{tabular}
              \caption{4-means cost.}
            \end{table}
        \end{enumerate}

      \item \label{qst:2B}
        \begin{enumerate}[label={(\roman*)}, ref={\theenumii(\roman*)}]
          \item
            The cumulative density function of the 4-means cost associated with the
            {\sf KMeans{\tt++}} algorithm is plotted below.
            \begin{figure}[H]
              \centering
              \begin{tikzpicture}
                \begin{axis}[
                  xlabel={4-Means Cost},
                  ylabel={Likelihood of Occurrence},
                ]
                  \addplot [
                    thick,
                    color=red,
                    hist={data=x, bins=100, handler/.style={const plot}, cumulative, density}
                  ] file {dev-dir/week4/data/kplusplus/kplusplus-cdf.csv};
                  \addlegendentry{CDF};
                \end{axis}
              \end{tikzpicture}
              \caption{
                4-Means Cost CDF Plot of {\sf KMeans{\tt++}} Algorithm.
              }
              \label{fig:4-means-cost-kmeans-cdf-plot}
            \end{figure}

          \item
            A total of $0 / 100 = 0.0$ of the trials, the subsets from {\sf Kmeans{\tt++}} and
            {\sf Gonzalez} matched within an error tolerance $\varepsilon < 0.01$. Please refer
            to \autoref{lst:kplusplus-cdf} for a fuller explanation on how this fractional value
            was obtained.
        \end{enumerate}

      \item \label{qst:2C}
        \begin{enumerate}[label={(\roman*)}, ref={\theenumii(\roman*)}]
          \item
            Below is the output of {\sf Lloyd}'s algorithm with centroids initially set as the
            points indexed $\{1,2,3,4\}$ in the {\sf C2.txt} dataset. A 4-means cost of
            approximately {\bf 0.8532158400114525} was obtained during the run.
            \begin{figure}[H]
              \centering
              \begin{tikzpicture}
                \begin{axis}[
                  legend pos={north east},
                  legend entries={Cluster 1, Cluster 2, Cluster 3, Cluster 4},
                  point meta={explicit},
                  mark size={1.5pt},
                  only marks,
                ]
                  \addplot [
                    scatter,
                    scatter/classes={
                      0={fill={cluster1}, draw={cluster1!80!black}},
                      1={fill={cluster2}, draw={cluster2!80!black}},
                      2={fill={cluster3}, draw={cluster3!80!black}},
                      3={fill={cluster4}, draw={cluster4!80!black}}
                    },
                  ] table[col sep=comma, x index=0, y index=1, meta index=2] {dev-dir/week4/data/lloyd/points-lloyd.csv};
                  \addplot [
                    mark=x,
                    mark size=5pt,
                    color=red
                  ] table[col sep=comma] {dev-dir/week4/data/lloyd/centers-lloyd.csv};
                \end{axis}
              \end{tikzpicture}
              \caption{
                {\sf Lloyd}'s algorithm with centroids initially set as the points indexed at
                $\{1,2,3,4\}$ on {\sf C2.txt} dataset. Please refer to \autoref{lst:lloyd} for
                a fuller explanation on how this plot was generated.
              }
              \label{fig:lloyd}
            \end{figure}

          \item
            Below is the output of {\sf Lloyd}'s algorithm with centroids initially set as the
            output {\sf Gonzalez}'s subroutine. The run incurred a 4-means cost of approximately
            {\bf 0.8769132227723297}.
            \begin{figure}[H]
              \centering
              \begin{tikzpicture}
                \begin{axis}[
                  legend pos={north east},
                  legend entries={Cluster 1, Cluster 2, Cluster 3, Cluster 4},
                  point meta={explicit},
                  mark size={1.5pt},
                  only marks,
                ]
                  \addplot [
                    scatter,
                    scatter/classes={
                      0={fill={cluster1}, draw={cluster1!80!black}},
                      1={fill={cluster2}, draw={cluster2!80!black}},
                      2={fill={cluster3}, draw={cluster3!80!black}},
                      3={fill={cluster4}, draw={cluster4!80!black}}
                    },
                  ] table[col sep=comma, x index=0, y index=1, meta index=2] {dev-dir/week4/data/lloyd/points-gonzalez.csv};
                  \addplot [
                    mark=x,
                    mark size=5pt,
                    color=red
                  ] table[col sep=comma] {dev-dir/week4/data/lloyd/centers-gonzalez.csv};
                \end{axis}
              \end{tikzpicture}
              \caption{
                {\sf Lloyd}'s algorithm with centroids initially set as the output of {\sf Gonzalez}'s
                subroutine. Please refer to \autoref{lst:lloyd} for a fuller explanation on how this
                plot has been generated.
              }
              \label{fig:lloyd-gonzalez}
            \end{figure}

          \item
            The cumulative density function of the 4-means cost associated with a vanilla
            implementation of {\sf Lloyd}'s algorithm (found in \autoref{lst:lloyd-cdf})
            is plotted below.
            \begin{figure}[H]
              \centering
              \begin{tikzpicture}
                \begin{axis}[
                  xlabel={4-Means Cost},
                  ylabel={Likelihood of Occurrence},
                ]
                  \addplot [
                    thick,
                    color=red,
                    hist={data=x, bins=100, handler/.style={const plot}, cumulative, density}
                  ] file {dev-dir/week4/data/lloyd/lloyd-cdf.csv};
                  \addlegendentry{CDF};
                \end{axis}
              \end{tikzpicture}
              \caption{
                4-Means Cost CDF Plot of {\sf Lloyd}'s Algorithm. (Vanilla Implementation).
              }
              \label{fig:4-means-cost-lloyd-kmeans-cdf-plot}
            \end{figure}
            A total of 0/100 = 0.0 of the trials, the subsets from {\sf Kmeans{\tt++}} and
            {\sf Lloyd} matched within an error tolerance $\varepsilon < 0.01$. Please refer
            to \autoref{lst:lloyd-cdf} for a fuller explanation on how this fractional value
            was obtained.
        \end{enumerate}

      \item \label{qst:2D}
        The cumulative density function of the 4-means cost associated with a {\sf scikit-learn}
        implementation of {\sf Lloyd}'s algorithm (found in \autoref{lst:lloyd-scikit-cdf})
        is plotted below.
        \begin{figure}[H]
          \centering
          \begin{tikzpicture}
            \begin{axis}[
              xlabel={4-Means Cost},
              ylabel={Likelihood of Occurrence},
            ]
              \addplot [
                thick,
                color=red,
                hist={data=x, bins=100, handler/.style={const plot}, cumulative, density}
              ] file {dev-dir/week4/data/lloyd/lloyd-scikit-cdf.csv};
              \addlegendentry{CDF};
            \end{axis}
          \end{tikzpicture}
          \caption{
            4-Means Cost CDF Plot of {\sf Lloyid}'s Algorithm. \\ (Scikit-Learn Implementation).
          }
          \label{fig:4-means-cost-lloyd-kmeans-cdf-plot}
        \end{figure}
        A total of 0/100 = 0.0 of the trials, the subsets from {\sf Kmeans{\tt++}} and {\sf Lloyd}
        matched within an error tolerance $\varepsilon < 0.01$. Please refer to
        \autoref{lst:lloyd-scikit-cdf} for a fuller explanation on how this fractional value was
        obtained.

        \lstinputlisting
          [caption={\sf lloyd-scikit-cdf.py}, label={lst:lloyd-scikit-cdf}]
          {dev-dir/week4/lloyd-scikit-cdf.py}
    \end{enumerate}

  \item \label{qst:3}
    Please refer to \autoref{lst:kmedians} for a fuller explanation on how the following centers
    have been generated. A cost of {\bf 0.4194209920788438} was obtained during the run. The
    associated {\sf .txt} file has been uploaded onto google drive under the alias {\tt\color{blue}
    u1419667-best-centers.txt}. \\ \bigskip

    \lstinputlisting
      [caption={\sf u1419667-best-centers.txt}, label={lst:u1419667-best-centers}]
      {dev-dir/week4/data/kmedians/best-centers.txt}
\end{enumerate}

\newpage

\section*{Appendix}

\lstinputlisting
  [caption={\sf hierarchical.py}, label={lst:hierarchical}]
  {dev-dir/week4/hierarchical.py}
  
\newpage

\lstinputlisting
  [caption={\sf hierarchical-dendrogram.py}, label={lst:hierarchical-dendrogram}]
  {dev-dir/week4/hierarchical-dendrogram.py}
  
\newpage

\lstinputlisting
  [caption={\sf gonzalez.py}, label={lst:gonzalez}]
  {dev-dir/week4/gonzalez.py}
  
\newpage

\lstinputlisting
  [caption={\sf kplusplus.py}, label={lst:kplusplus}]
  {dev-dir/week4/kplusplus.py}
  
\newpage

\lstinputlisting
  [caption={\sf kplusplus-cdf.py}, label={lst:kplusplus-cdf}]
  {dev-dir/week4/kplusplus-cdf.py}
  
\newpage

\lstinputlisting
  [caption={\sf lloyd.py}, label={lst:lloyd}]
  {dev-dir/week4/lloyd.py}
  
\newpage

\lstinputlisting
  [caption={\sf lloyd-cdf.py}, label={lst:lloyd-cdf}]
  {dev-dir/week4/lloyd-cdf.py}
  
\newpage

\lstinputlisting
  [caption={\sf kmedians.py}, label={lst:kmedians}]
  {dev-dir/week4/kmedians.py}
\end{document}
